\section{Concetto di Hydration}
Questa tecnica consiste nel servire pagine statiche attraverso static hosting o tramite SSR (Server Side Rendering) e
poi, una volta che la pagina è stata caricata, aggiungere interattività attraverso JavaScript attaccando \textit{event
handlers} agli elementi HTML.
Grazie a questa tecnica, la pagina viene renderizzata sul server e inviata al client con del contenuto che, nonostate
sia statico, la fa apparire pronta all'uso e prensenta una porzione dei dati all'utente, dando l'apparenza di un
caricamento istantaneo e di maggiore fluidità.